% Problem-section file following ../_template.tex.
\section{Capacity-achieving codes for amplitude damping}
\label{sec:problem-2}

\paragraph{Problem.}
What is the constructive quantum code for achieving the quantum capacity
$\mathcal{Q}(\mathcal A_p)$ of the qubit amplitude-damping channel
\begin{equation}
  \mathcal A_p(\rho)=A_0\rho A_0^\dagger+A_1\rho A_1^\dagger,
  \qquad 0\le p\le1?
  \label{eq:p2-amplitude-damping}
\end{equation}
The operators appearing in Eq.~\eqref{eq:p2-amplitude-damping} are
\begin{equation}
  \begin{aligned}
    A_0&=\lvert0\rangle\!\langle0\rvert
         +\sqrt{1-p}\,\lvert1\rangle\!\langle1\rvert
       =\begin{pmatrix}1&0\\0&\sqrt{1-p}\end{pmatrix},\\
    A_1&=\sqrt p\,\lvert0\rangle\!\langle1\rvert
       =\begin{pmatrix}0&\sqrt p\\0&0\end{pmatrix}.
  \end{aligned}
  \label{eq:p2-kraus-operators}
\end{equation}
Equation~\eqref{eq:p2-kraus-operators} uses $p$ as the decay probability of
the excited state.

\subsection*{Status}

Unsolved

\subsection*{Source}

This constructive gap is implicit in the exact capacity formula of Wolf and
P\'erez-Garc\'ia and the polar-code construction of Wilde and Guha, which
attains only the symmetric coherent-information rate for this channel
\sourcecite{ref:p2-wolf}{WPG07},
\sourcecite{ref:p2-wilde-guha}{WG13}.

\subsection*{Progress}

\begin{itemize}
  \item The channel is degradable for $0\le p\le1/2$ and antidegradable for
  $1/2\le p\le1$.  Consequently,
  \begin{equation}
    \mathcal{Q}(\mathcal A_p)=
    \begin{cases}
      \displaystyle\max_{0\le q\le1}
      \bigl[h_2((1-p)q)-h_2(pq)\bigr],&0\le p\le\tfrac12,\\[1.5mm]
      0,&\tfrac12\le p\le1,
    \end{cases}
    \label{eq:p2-capacity}
  \end{equation}
  where $q$ is the excited-state population of the diagonal input and
  $h_2(x):=-x\log_2x-(1-x)\log_2(1-x)$, with $0\log_2 0:=0$.  The capacity formula in
  Eq.~\eqref{eq:p2-capacity} follows from the small-environment analysis
  \sourcecite{ref:p2-wolf}{WPG07}.

  \item For $0\le p\le1/2$, Wilde and Guha constructed a channel-adapted
  quantum polar code for degradable channels with a classical environment,
  explicitly including the amplitude-damping channel.  It achieves
  \begin{equation}
    I_c(I/2,\mathcal A_p)
    =h_2\!\left(\frac{1-p}{2}\right)-h_2\!\left(\frac p2\right),
    \label{eq:p2-symmetric-rate}
  \end{equation}
  the symmetric coherent-information rate.  The rate in
  Eq.~\eqref{eq:p2-symmetric-rate} is obtained with an efficient encoder and
  asymptotically vanishing entanglement consumption, although that work did
  not establish an efficient decoder \sourcecite{ref:p2-wilde-guha}{WG13}.
  Wilde and Renes subsequently gave a polar construction for arbitrary
  qubit-input channels that achieves the same rate using coherent
  successive-cancellation decoding \sourcecite{ref:p2-wilde-renes}{WR12}.
\end{itemize}

\subsection*{References}

\begin{enumerate}[label={},leftmargin=4.8em,labelsep=0.5em]
  \footnotesize
  \item[\textup{[WPG07]}]\label{ref:p2-wolf}
  M. M. Wolf and D. P\'erez-Garc\'ia, ``Quantum Capacities of Channels with
  Small Environment,'' \emph{Physical Review A} \textbf{75}, 012303 (2007).
  \href{https://doi.org/10.1103/PhysRevA.75.012303}{doi:10.1103/PhysRevA.75.012303};
  \href{https://arxiv.org/abs/quant-ph/0607070}{arXiv:quant-ph/0607070}.

  \item[\textup{[WG13]}]\label{ref:p2-wilde-guha}
  M. M. Wilde and S. Guha, ``Polar Codes for Degradable Quantum Channels,''
  \emph{IEEE Transactions on Information Theory} \textbf{59}, 4718--4729
  (2013). \href{https://doi.org/10.1109/TIT.2013.2250575}{doi:10.1109/TIT.2013.2250575};
  \href{https://arxiv.org/abs/1109.5346}{arXiv:1109.5346}.

  \item[\textup{[WR12]}]\label{ref:p2-wilde-renes}
  M. M. Wilde and J. M. Renes, ``Quantum Polar Codes for Arbitrary Channels,''
  in \emph{2012 IEEE International Symposium on Information Theory
  Proceedings}, 334--338 (2012).
  \href{https://doi.org/10.1109/ISIT.2012.6284203}{doi:10.1109/ISIT.2012.6284203};
  \href{https://arxiv.org/abs/1201.2906}{arXiv:1201.2906}.
\end{enumerate}

\subsection*{Comment}

Because the amplitude-damping channel is nonunital, the maximizing population
$q$ in Eq.~\eqref{eq:p2-capacity} is generally not $1/2$.  A code achieving
only Eq.~\eqref{eq:p2-symmetric-rate} therefore does not by itself achieve the
full capacity.  The unresolved constructive task is an explicit, efficient
capacity-achieving scheme incorporating the required asymmetric input shaping.

\subsection*{Tag}

Amplitude-damping channels; Quantum capacity; Quantum coding theory;
Quantum polar codes; Degradable channels; Qubit systems

\subsection*{ID}

\texttt{op\_fcd21a1a5021e464}
